\documentclass[10pt]{article}

\usepackage[margin=1in]{geometry}
\usepackage{amsmath,amssymb,amsthm,mathtools}
\usepackage{booktabs}
\usepackage{graphicx}
\usepackage{microtype}
\usepackage{xcolor}
\usepackage{pgfplots}
\usepackage{hyperref}

\pgfplotsset{compat=1.18}

\newtheorem{theorem}{Theorem}
\newtheorem{proposition}{Proposition}
\newtheorem{lemma}{Lemma}
\newtheorem{remark}{Remark}

\newcommand{\one}{\mathbf{1}}
\newcommand{\R}{\mathbb{R}}
\newcommand{\Birk}{\mathcal{B}_n}
\newcommand{\Mone}{\mathcal{M}_{\one}}
\newcommand{\pperp}{P_{\perp}}
\newcommand{\operp}{\one^{\perp}}
\newcommand{\Eperp}{E_{\perp}}
\newcommand{\sv}{\sigma}
\DeclareMathOperator{\diag}{diag}
\DeclareMathOperator{\polar}{polar}
\DeclareMathOperator{\rank}{rank}

\title{
Birkhoff Stability Is Not Isometric Stability:\\
Fixed-Vector Residual Transport for Hyper-Connections
}

\author{
Author information to be finalized
}

\date{May 2026}

\begin{document}
\maketitle

\begin{abstract}
Hyper-Connections expand residual connections into multiple residual streams,
making residual transport itself a learnable geometric object.  This creates a
final-mile problem for residual architectures: unconstrained multi-stream
transport is expressive but can destroy the identity mapping and amplify signals
through depth, while Birkhoff-constrained transport, as used in
manifold-constrained Hyper-Connections (mHC), restores the residual mean but can
dissipate the mean-zero stream complement.  We study this gap and propose
IsoHC, a fixed-vector isometric residual transport constraint
$H^\top H=I,\ H\one=\one$.  IsoHC preserves the residual mean exactly while
keeping every singular value on $\operp$ equal to one.  We show theoretically
that mean preservation does not imply complement preservation: doubly stochastic
transport can be perfectly stable in the residual mean direction while
contracting $\operp$ exponentially through depth.  In residual-only detectors,
unconstrained transport explodes, Birkhoff-like diffusion collapses the
mean-zero energy, and IsoHC preserves gradients, energy, and stream diversity
through 1024 layers.  In a 48-layer FineWeb-Edu Transformer posthoc analysis,
mHC's composite mean-zero transport gain falls to $4.39\times10^{-4}$, whereas
IsoHC remains near $0.996$.  At the current compute scale, we do not claim a
decisive language-modeling perplexity improvement.  Instead, we identify and
repair a concrete residual-transport pathology: preserving the residual mean is
not sufficient for preserving multi-stream residual information.
\end{abstract}

\section{Introduction}

Residual connections stabilize deep neural networks by preserving an identity
path.  A standard residual block can be abstracted as
\begin{equation}
    x_{l+1}=x_l+F_l(x_l),
\end{equation}
so that each layer writes an increment into a persistent residual state.
Pre-normalized Transformers inherit this design: attention and MLP branches read
from a shared residual stream and add updates back into it.  This identity path
is one of the main reasons very deep networks can be optimized.

Hyper-Connections (HC) generalize this paradigm by replacing the single
residual stream with multiple residual streams and a learnable transport between
them~\cite{hyperconnections}.  This turns the residual pathway from a scalar
identity channel into a structured multi-stream transport system.  However, it
also makes the residual connection itself a geometric object.  A generic learned
transport matrix may amplify, collapse, or drift away from the residual identity
semantics, and products of such matrices can become unstable through depth.

mHC repairs this instability by projecting residual mixing matrices onto the
Birkhoff polytope~\cite{mhc}:
\begin{equation}
    \Birk=\{H\in\R^{n\times n}:H\one=\one,\ \one^\top H=\one^\top,\ H\ge 0\}.
\end{equation}
This is a natural stabilization: doubly stochastic transport preserves the
uniform stream mean, restoring the identity-mapping property that residual
connections are meant to provide.  Recent work has further improved this line
through alternative parameterizations and spectral constraints
~\cite{shc,kromhc}.

Our observation is that mean preservation is not the same as residual-complement
preservation.  Let
\begin{equation}
    \pperp = I - \frac{1}{n}\one\one^\top
\end{equation}
be the projection onto the mean-zero stream subspace $\operp$.  In a
multi-stream residual architecture, the stream mean carries the ordinary
identity path, while $\pperp X$ carries the differences between streams.  If
these differences are intended to store persistent residual information, then a
residual transport geometry should not only preserve the mean; it should also
avoid dissipating the complement.

Birkhoff geometry does not guarantee this.  A doubly stochastic matrix is
non-expansive, but it is generally diffusive.  It can make all standard
mean-direction stability diagnostics look healthy while exponentially
contracting $\operp$.  We call this distinction:
\begin{center}
\emph{Birkhoff stability is not isometric stability.}
\end{center}

We propose IsoHC, a fixed-vector isometric constraint for residual transport:
\begin{equation}
    \Mone = \{Q\in O(n):Q\one=\one\}.
\end{equation}
IsoHC preserves the residual mean exactly and acts orthogonally on
$\operp$.  It therefore avoids both unconstrained amplification and
Birkhoff-style diffusion.  The contribution of this note is not a claim of
large-scale language-modeling superiority.  Rather, it is a mechanism diagnosis
and repair: mHC protects the residual mean, but it can dissipate the residual
complement; fixed-vector isometry protects both.

\paragraph{Contributions.}
We make four contributions.
\begin{enumerate}
    \item We formulate the residual-complement gap in multi-stream residual
    transport: preserving the stream mean does not preserve the full residual
    transport geometry.
    \item We prove that Birkhoff constraints can hide depth-wise contraction on
    $\operp$, even when mean-direction stability is exact.
    \item We derive IsoHC as the fixed-vector orthogonal stabilizer that
    preserves both the residual mean and the complement norm.
    \item We provide controlled evidence from residual-only detectors, graph
    oversmoothing experiments, and 48-layer Transformer posthoc analysis.
\end{enumerate}

\section{Multi-Stream Residual Transport}

Let $X_l\in\R^{n\times d}$ denote $n$ residual streams at layer $l$.  A generic
multi-stream residual update can be written as
\begin{equation}
    X_{l+1}=H_lX_l+b_l\otimes y_l,
\end{equation}
where $H_l\in\R^{n\times n}$ transports the streams and $y_l\in\R^d$ is the
branch output.  The stream mean is
\begin{equation}
    \mu(X)=\frac{1}{n}\one^\top X.
\end{equation}
To embed the standard residual path into the mean stream, we want
\begin{equation}
    \mu(X_{l+1})=\mu(X_l)+y_l.
\end{equation}
This is satisfied if
\begin{equation}
    \one^\top H_l=\one^\top,\qquad \frac{1}{n}\one^\top b_l=1.
\end{equation}

Every stream state decomposes as
\begin{equation}
    X=\one\otimes\mu(X)+\pperp X.
\end{equation}
The first term is the residual mean.  The second term is the residual
complement, where the additional degrees of freedom of a multi-stream residual
architecture live.  This gives two distinct requirements:
\begin{align}
    \text{mean direction:} &\quad \text{preserve the residual identity path},\\
    \text{mean-zero complement:} &\quad \text{transport stream differences without
    expansion or dissipation}.
\end{align}
mHC solves the first requirement by using Birkhoff constraints.  IsoHC targets
both.

\section{Birkhoff Stability Versus Complement Preservation}

Let $U\in\R^{n\times(n-1)}$ be an orthonormal basis of $\operp$, so that
$U^\top U=I$ and $U^\top\one=0$.  The complement action of a mean-preserving
matrix $H$ is
\begin{equation}
    H_\perp = U^\top H U.
\end{equation}

\begin{theorem}[Birkhoff transport is non-expansive, not necessarily isometric]
\label{thm:birkhoff-nonexpansive}
If $H\in\Birk$, then $H$ preserves the stream mean and
\begin{equation}
    \|U^\top H U\|_2 \le 1.
\end{equation}
However, Birkhoff constraints do not imply that $U^\top H U$ has singular
values equal to one.
\end{theorem}

\begin{proof}
By the Birkhoff-von Neumann theorem, $H$ is a convex combination of permutation
matrices:
\begin{equation}
    H=\sum_k \alpha_k P_k,\qquad \alpha_k\ge 0,\quad \sum_k\alpha_k=1.
\end{equation}
Each $P_k$ is orthogonal, hence $\|P_k\|_2=1$.  Therefore
\begin{equation}
    \|H\|_2 \le \sum_k\alpha_k\|P_k\|_2 = 1.
\end{equation}
Since $H\one=\one$ and $\one^\top H=\one^\top$, the subspace $\operp$ is
invariant under $H$, and the restriction $U^\top H U$ is also non-expansive.
Non-expansiveness does not force all singular values to equal one.
\end{proof}

\begin{theorem}[Mean stability can hide exponential complement contraction]
\label{thm:halpha}
For $0\le \alpha<1$, define
\begin{equation}
    H_\alpha = \alpha I + (1-\alpha)\frac{\one\one^\top}{n}.
\end{equation}
Then $H_\alpha\in\Birk$, $H_\alpha\one=\one$, and for any
$x\in\operp$,
\begin{equation}
    H_\alpha x = \alpha x.
\end{equation}
Consequently,
\begin{equation}
    \|\pperp H_\alpha^L\pperp x\|_2
    =
    \alpha^L \|\pperp x\|_2.
\end{equation}
\end{theorem}

\begin{proof}
$H_\alpha$ is nonnegative for $0\le\alpha<1$, and its rows and columns sum to
one, so $H_\alpha\in\Birk$.  If $x\in\operp$, then $\one^\top x=0$ and
$\one\one^\top x/n=0$.  Thus $H_\alpha x=\alpha x$, and the depth-$L$ statement
follows by repeated application.
\end{proof}

Theorem~\ref{thm:halpha} is the core counterexample: row/column stochasticity can
be perfect while complement information decays exponentially.  This is not a
bug in mHC.  mHC is designed to restore identity stability.  The point is that
identity stability alone does not certify multi-stream complement preservation.

\section{IsoHC: Fixed-Vector Isometric Transport}

IsoHC replaces Birkhoff residual transport with the fixed-vector orthogonal
stabilizer
\begin{equation}
    \Mone=\{Q\in O(n):Q\one=\one\}.
\end{equation}

\begin{theorem}[Fixed-vector isometry preserves mean and complement norm]
\label{thm:isohc}
Let $Q\in\Mone$.  Then for any $X\in\R^{n\times d}$,
\begin{equation}
    \mu(QX)=\mu(X),
\end{equation}
and
\begin{equation}
    \|\pperp QX\|_F=\|\pperp X\|_F.
\end{equation}
Equivalently, all singular values of $U^\top Q U$ are one.
\end{theorem}

\begin{proof}
Since $Q^\top Q=I$ and $Q\one=\one$, we also have $Q^\top\one=\one$.  Hence
\begin{equation}
    \mu(QX)=\frac{1}{n}\one^\top QX=\frac{1}{n}\one^\top X=\mu(X).
\end{equation}
The subspace $\operp$ is invariant under $Q$, and the restriction
$U^\top Q U$ is orthogonal.  Therefore it preserves the Frobenius norm of the
complement.
\end{proof}

\begin{proposition}[Nonnegative orthogonal stochastic matrices are discrete]
\label{prop:permutation}
If $Q\ge 0$, $Q\one=\one$, and $Q^\top Q=I$, then $Q$ is a permutation matrix.
\end{proposition}

\begin{proof}
The rows of $Q$ are nonnegative unit vectors with pairwise inner product zero.
Thus distinct rows have disjoint support.  Since each row sums to one and has
unit Euclidean norm, each row has exactly one nonzero entry, equal to one.
Column orthogonality then implies that each column has exactly one such entry.
\end{proof}

This proposition explains why continuous non-diffusive residual transport must
leave the nonnegative Birkhoff geometry.  If we require nonnegativity,
stochasticity, and orthogonality simultaneously, the feasible set degenerates to
permutations.

\subsection{Projection and finite-step retraction}

Let $v=\one/\sqrt{n}$ and let $U$ span $v^\perp$.  Every element of $\Mone$ has
the form
\begin{equation}
    Q = vv^\top + U R U^\top,\qquad R\in O(n-1).
\end{equation}
Given an unconstrained matrix $A$, the Frobenius-nearest exact projection, away
from polar degeneracies, is
\begin{equation}
    \Pi_v(A)=vv^\top + U\,\polar(U^\top A U)\,U^\top.
\end{equation}
In implementation, we use a finite-step Newton-Schulz polar iteration as an
efficient fixed-vector retraction:
\begin{equation}
    B=U^\top A U,\qquad R_K\approx \polar(B),\qquad Q_K=vv^\top+UR_KU^\top.
\end{equation}
The finite-step retraction preserves the fixed vector by construction, but
orthogonality is approximate and must be monitored.  In our 48-layer run,
IsoHC's fixed-vector error is $2.79\times10^{-7}$ and its orthogonality error is
$5.57\times10^{-4}$.

\section{Experimental Evidence}

The experiments are designed to test residual transport geometry, not to claim
language-modeling state of the art.  The central question is:
\begin{quote}
Does the transport preserve both the residual mean and the residual complement
through depth?
\end{quote}

\subsection{Residual-only depth detector}

We first remove attention and MLP branches and study pure transport through up
to 1024 layers.  Table~\ref{tab:detector} shows the $n=8$, $L=1024$ detector.

\begin{table}[t]
\centering
\caption{Residual-only transport detector at $n=8$, $L=1024$.  IsoHC preserves
gradient scale, mean-zero energy, and stream diversity.  Diffusive transport
collapses the complement, while unconstrained transport explodes.}
\label{tab:detector}
\begin{tabular}{lrrrrr}
\toprule
method & grad mean & mean-zero energy & stream cosine & $\sigma_{\min}$ & $\sigma_{\max}$ \\
\midrule
IsoHC & 1.0000 & 1.0000 & 0.0250 & 1.0000 & 1.0000 \\
mHC-lite & 0.3536 & 0.0000 & 1.0000 & 0.5724 & 1.0000 \\
GCN diffusion & 0.3533 & 0.0000 & 1.0000 & 0.8000 & 1.0000 \\
unconstrained & $5.58{\times}10^{13}$ & $5.61{\times}10^{13}$ & 0.9877 & 0.4270 & 1.6155 \\
\bottomrule
\end{tabular}
\end{table}

The detector isolates the geometry.  Unconstrained HC can amplify through depth;
Birkhoff-like diffusion prevents explosion but collapses the complement; IsoHC
preserves the complement isometrically.

\subsection{48-layer FineWeb-Edu mechanism run}

We trained a 48-layer GPT-style Transformer on a FineWeb-Edu token cache:
\begin{equation}
    L=48,\quad d_{\mathrm{model}}=512,\quad h=8,\quad n=4,\quad T=512,
\end{equation}
for 20M tokens per method.  The model has roughly 177M parameters.  A fair batch
probe selected a common batch size across methods.

\begin{table}[t]
\centering
\caption{48-layer FineWeb-Edu short-budget run.  Perplexity differences are not
the main claim.  The key observation is that mHC reduces $\Eperp$, while IsoHC
preserves much larger complement energy.}
\label{tab:fe48}
\begin{tabular}{lrrrr}
\toprule
method & val loss & val PPL & $\Eperp$ init $\to$ final & stream cosine init $\to$ final \\
\midrule
baseline & 5.7366 & 310.01 & -- & -- \\
identity-HC & 5.7157 & 303.60 & 0.0222 $\to$ 0.0503 & 0.9993 $\to$ 0.9971 \\
unconstrained HC & 5.6760 & 291.78 & 0.0222 $\to$ 0.0547 & 0.9993 $\to$ 1.0000 \\
mHC & 5.7249 & 306.39 & 0.0224 $\to$ 0.0136 & 0.9993 $\to$ 0.9998 \\
IsoHC & 5.7206 & 305.10 & 0.0224 $\to$ 0.0546 & 0.9993 $\to$ 0.9967 \\
\bottomrule
\end{tabular}
\end{table}

\begin{table}[t]
\centering
\caption{Single-layer transport diagnostics in the 48-layer run.  mHC satisfies
stochastic constraints but has average complement singular value near 0.922.
IsoHC keeps the fixed vector and complement singular values near one.}
\label{tab:diagnostics}
\small
\resizebox{\linewidth}{!}{%
\begin{tabular}{llll}
\toprule
method & invariant error & orthogonality error & $\operp$ singular values \\
\midrule
identity-HC & 0.0 & 0.0 & exact identity \\
unconstrained HC & 0.1050 & 0.1694 & unconstrained \\
mHC & row 0.0046 / col $2.30{\times}10^{-7}$ & -- & mean 0.9222, min 0.9208, max 0.9237 \\
IsoHC & $2.79{\times}10^{-7}$ & $5.57{\times}10^{-4}$ & mean 0.99996, min 0.99982, max 1.00012 \\
\bottomrule
\end{tabular}
}
\end{table}

\begin{figure}[t]
\centering
\begin{tikzpicture}
\begin{axis}[
    width=0.92\linewidth,
    height=0.38\linewidth,
    xlabel={transport index},
    ylabel={single-layer mean $\sigma(U^\top H_l U)$},
    ymin=0.90,
    ymax=1.02,
    legend pos=south east,
    grid=both,
    grid style={gray!20},
]
\addplot[blue, thick] table[x=index,y=sv_mean,col sep=comma] {data/mhc_single_spectrum.csv};
\addlegendentry{mHC}
\addplot[red, thick] table[x=index,y=sv_mean,col sep=comma] {data/isohc_single_spectrum.csv};
\addlegendentry{IsoHC}
\end{axis}
\end{tikzpicture}
\caption{Single-layer complement singular values in the 48-layer Transformer.
mHC is consistently contractive on $\operp$, while IsoHC remains near one.}
\label{fig:single-spectrum}
\end{figure}

\begin{figure}[t]
\centering
\begin{tikzpicture}
\begin{semilogyaxis}[
    width=0.92\linewidth,
    height=0.42\linewidth,
    xlabel={transport index},
    ylabel={composite mean $\sigma(U^\top H_{1:l}U)$},
    ymin=1e-4,
    ymax=2,
    legend pos=south west,
    grid=both,
    grid style={gray!20},
]
\addplot[blue, thick] table[x=index,y=composite_sv_mean,col sep=comma] {data/mhc_composite_gain.csv};
\addlegendentry{mHC}
\addplot[red, thick] table[x=index,y=composite_sv_mean,col sep=comma] {data/isohc_composite_gain.csv};
\addlegendentry{IsoHC}
\end{semilogyaxis}
\end{tikzpicture}
\caption{Composite complement transport gain.  At the final transport,
mHC falls to $4.39\times 10^{-4}$, whereas IsoHC remains near $0.996$.  This
is the central mechanism result: mean-preserving stability does not certify
complement-preserving transport.}
\label{fig:composite}
\end{figure}

\subsection{Checkpoint posthoc intervention}

We further analyzed checkpointed mHC and IsoHC runs using 16 validation batches
with batch size 1 for intervention and gradient probes.

\begin{table}[t]
\centering
\caption{48-layer checkpoint posthoc analysis.  The transport pathology is
strong, but functional loss deltas are small at this budget.}
\label{tab:posthoc}
\begin{tabular}{lrrrrr}
\toprule
method & val loss & PPL & final composite $\operp$ gain & best removal $\Delta L$ & replace by identity $\Delta L$ \\
\midrule
mHC & 5.900727 & 365.30 & 0.000439 & 0.000289 & -0.000524 \\
IsoHC & 5.901274 & 365.50 & 0.996349 & 0.000516 & 0.000713 \\
\bottomrule
\end{tabular}
\end{table}

The posthoc result is the cleanest Transformer mechanism evidence: mHC's
composite complement transport nearly vanishes, while IsoHC remains isometric.
However, complement-removal and IsoHC-to-identity replacement change average
validation loss only at the $10^{-4}$ to $10^{-3}$ scale.  This means the current
run establishes the transport pathology, but does not yet show that the
preserved complement is strongly used by the language model.

\subsection{Graph oversmoothing side evidence}

Graph neural networks provide an additional stress test for diffusion.  In a
synthetic oversmoothing experiment at depth 128, ordinary diffusion collapses
node variation, while an isometric transport preserves it:

\begin{table}[t]
\centering
\caption{Synthetic graph oversmoothing stress test at depth 128.}
\label{tab:gnn}
\begin{tabular}{lrrrr}
\toprule
method & energy & centered variance & cosine & Dirichlet energy \\
\midrule
GCN & 0.2003 & 0.0012 & 1.0000 & 0.0000 \\
Residual GCN & 0.2003 & 0.0012 & 1.0000 & 0.0000 \\
IsoNode & 1.0000 & 1.0000 & 0.0325 & 7940.2 \\
\bottomrule
\end{tabular}
\end{table}

On Cora node classification, a corrected stream architecture recovers deep
performance where ordinary GCNs collapse:
\begin{table}[t]
\centering
\caption{Cora deep node classification.  This is side evidence for the
anti-collapse role of isometric stream transport, not the main LM claim.}
\label{tab:cora}
\begin{tabular}{lccc}
\toprule
model & L2 & L16 & L32 \\
\midrule
GCN & $77.6\pm1.9$ & $31.1\pm0.2$ & $31.2\pm0.3$ \\
ResGCN & $75.8\pm1.2$ & $28.8\pm4.0$ & $28.6\pm2.8$ \\
IsoStream v2 & -- & $73.1\pm1.9$ & $66.6\pm3.4$ \\
\bottomrule
\end{tabular}
\end{table}

\section{Discussion}

\paragraph{What IsoHC solves.}
IsoHC should not be read as a small-model perplexity trick.  It targets residual
transport geometry.  Ordinary residual connections provide a stable identity
path but only one residual transport channel.  HC adds multiple streams, but
unconstrained transport can destabilize the identity path.  mHC restores the
mean identity path with Birkhoff geometry.  IsoHC completes this geometry by
also preserving the residual complement.

\paragraph{Why not just identity-HC?}
The identity matrix is a member of $\Mone$ and preserves both mean and
complement.  It is therefore an essential control.  IsoHC's additional value is
learnable non-destructive stream communication: $U^\top H_lU$ can be a nontrivial
orthogonal transform rather than the identity.  Our current LM replacement probe
shows only a small IsoHC-to-identity loss delta, so the present evidence verifies
the geometry but does not yet demonstrate strong functional value of learned
rotations over identity-HC.

\paragraph{Relation to spectral HC.}
Spectral constraints such as $\|H\|_2\le 1$ prevent expansion but do not by
themselves guarantee residual mean preservation or complement isometry.  IsoHC
imposes the stronger geometry
\begin{equation}
    H\one=\one,\qquad H^\top H=I,
\end{equation}
which simultaneously preserves the residual mean and all complement singular
values.  Thus the distinction is not simply signed versus nonnegative mixing,
but non-expansive versus isometric transport on the invariant complement.

\section{Limitations}

This note intentionally makes a limited claim.
\begin{enumerate}
    \item We do not show decisive large-scale LM perplexity improvements over
    mHC or spectral HC variants.
    \item The current LM complement-intervention effects are small, so
    functional use of the preserved complement remains weak at this scale.
    \item Identity-HC remains a strong control.  Future work must show when
    learned isometric stream mixing is more useful than no mixing.
    \item IsoHC gives fixed-vector isometry for the residual transport
    component, not full dynamical isometry of an entire Transformer block.
    \item Finite-step Newton-Schulz retractions are approximate; fixed-vector
    error and orthogonality error must be monitored.
\end{enumerate}

\section{Future Work}

The next experiments should test whether the preserved complement is
functionally used, rather than merely energetically preserved.

\paragraph{Token-level intervention.}
For each token, compare full logits $p_t$ with intervention logits $p'_t$ after
complement removal or complement clamping:
\begin{equation}
    D_{\mathrm{KL}}(p_t\|p'_t),\qquad
    \Delta\ell_t = -\log p'_t(y_t)+\log p_t(y_t).
\end{equation}
The average loss delta may be small even if high-entropy, rare, late-position,
or copy-sensitive tokens depend on the complement.

\paragraph{Complement clamp and gate sweeps.}
Instead of removing the complement at one layer, clamp it from layer $k$ onward:
\begin{equation}
    X_l\leftarrow \one\mu(X_l),\qquad l\ge k,
\end{equation}
or sweep
\begin{equation}
    X_l^{(\alpha)}=\one\mu(X_l)+\alpha\pperp X_l,\qquad
    \alpha\in\{0,0.25,0.5,0.75,1\}.
\end{equation}
Monotonic sensitivity to $\alpha$ would provide a stronger causal bridge.

\paragraph{Deep-thin stress and routing tasks.}
Depth is the natural axis for a contraction claim.  Future work should run
48/72/96/128-layer deep-thin models and synthetic routing or associative recall
tasks that require cross-stream communication.  These settings can test whether
learned isometric mixing is functionally necessary and whether mHC's complement
contraction limits behavior.

\section{Conclusion}

Hyper-Connections make residual transport multi-stream.  mHC restores the
residual mean by imposing Birkhoff constraints, but preserving the mean is not
the same as preserving the full residual transport geometry.  Birkhoff transport
is non-expansive and can be strongly contractive on the mean-zero complement.
IsoHC imposes fixed-vector isometry, preserving both the residual mean and the
complement norm.  The current evidence supports IsoHC as a geometric mechanism:
mHC's composite complement transport collapses in depth, while IsoHC remains
approximately isometric.  The remaining open question is functional relevance:
when, and at what scale, does the preserved residual complement become necessary
for model behavior?

\begin{thebibliography}{9}

\bibitem{resnet}
K. He, X. Zhang, S. Ren, and J. Sun.
\newblock Deep Residual Learning for Image Recognition.
\newblock In \emph{CVPR}, 2016.

\bibitem{transformer}
A. Vaswani et al.
\newblock Attention Is All You Need.
\newblock In \emph{NeurIPS}, 2017.

\bibitem{hyperconnections}
D. Zhu, H. Huang, Z. Huang, Y. Zeng, Y. Mao, B. Wu, Q. Min, and X. Zhou.
\newblock Hyper-Connections.
\newblock arXiv:2409.19606, 2024.

\bibitem{mhc}
Z. Xie et al.
\newblock mHC: Manifold-Constrained Hyper-Connections.
\newblock arXiv:2512.24880, 2025.

\bibitem{shc}
Z. Liu, H. Zhang, and A. Li.
\newblock Beyond the Birkhoff Polytope: Spectral-Sphere-Constrained Hyper-Connections.
\newblock arXiv:2603.20896, 2026.

\bibitem{kromhc}
W. Zhou, Y. Gu, G. Iacovides, and D. Mandic.
\newblock KromHC: Manifold-Constrained Hyper-Connections with Kronecker-Product Residual Matrices.
\newblock arXiv:2601.21579, 2026.

\bibitem{pairnorm}
L. Zhao and L. Akoglu.
\newblock PairNorm: Tackling Oversmoothing in GNNs.
\newblock In \emph{ICLR}, 2020.

\end{thebibliography}

\end{document}
